% Supplementary Material for StreamMind (separate PDF)
\documentclass[10pt,twocolumn,letterpaper]{article}

\usepackage{cvpr}

\usepackage[accsupp]{axessibility}
\usepackage{booktabs}
\usepackage{graphicx}
\usepackage{amsmath,amssymb}
\usepackage{xspace}
\usepackage{microtype}

\definecolor{cvprblue}{rgb}{0.21,0.49,0.74}
\usepackage[pagebackref,breaklinks,colorlinks,allcolors=cvprblue]{hyperref}

\newcommand{\ours}{StreamMind\xspace}

\def\paperID{*****}
\def\confName{CVPR}
\def\confYear{2026}

\title{Adaptive Temporal Memory for Interactive \\ Question Answering on Live Video Streams \\ {\large Supplementary Material}}

\author{
Suresh Kumar Palus\textsuperscript{1} \qquad Partha Sarathi Samal\textsuperscript{2}\\[4pt]
\textsuperscript{1}Independent Researcher, Pennsylvania, USA\\
\textsuperscript{2}Independent Researcher, Connecticut, USA\\[2pt]
{\tt\small \textsuperscript{1}sure.1985@gmail.com \quad \textsuperscript{2}samalpartha@gmail.com}
}

\begin{document}
\maketitle

\section{TQR Keyword Lists}
\label{sec:suppl_tqr}

\Cref{tab:tqr_keywords} lists the full keyword sets used by the Temporal Query Router. Keywords were compiled from annotated question-scope pairs in the LiveQA-Bench training split and validated on a held-out split.

\begin{table}[h]
  \caption{\textbf{TQR keyword lists.} Each list is matched against the lowercased query string; the scope with the most hits wins.}
  \label{tab:tqr_keywords}
  \centering
  \small
  \begin{tabular}{@{}lp{5.8cm}@{}}
    \toprule
    Scope & Keywords \\
    \midrule
    \textsc{Instant} & right now, currently, at this moment, happening now, see right now \\[3pt]
    \textsc{Recent} & just, a moment ago, few seconds, recently, did the, was the, just happened, last minute, a minute ago \\[3pt]
    \textsc{Historical} & earlier, before, previously, at the start, in the beginning, how many times, ever, at any point, throughout, so far \\
    \bottomrule
  \end{tabular}
\end{table}

\section{LiveQA-Bench Construction Details}
\label{sec:suppl_liveqa}

LiveQA-Bench uses 55 video streams covering diverse visual domains:
\emph{daily life/indoor} (cooking, surveillance, gym, grocery, classroom, warehouse, caf\'{e}, workshop, library, restaurant),
\emph{outdoor/urban} (traffic, street, aerial city, night city, market, city jogger),
\emph{nature} (beach, forest, sunset, rain, underwater, forest walk),
\emph{sports/action} (boxing, cycling, running, swimming, soccer, basketball, MMA, kickboxing, surfing, motocross, skateboarding, skiing),
and \emph{cinematic/trailer} (car chase, concert, fireworks, helicopter rescue, military ops, tunnel crash, sports cars, dashboard driving).
The activity composite concatenates four distinct indoor/outdoor scenes.
Source clips come from Ego4D~\cite{grauman2022ego4d} and free-licensed Mixkit videos.
Each stream is 10--80 seconds long.
Ground-truth answers are generated automatically: for \textsc{instant} questions, BLIP captions and targeted VQA probes (location, objects) of the target frame serve as ground truth; for \textsc{recent} and \textsc{historical} questions, Flan-T5 summarizes the BLIP captions from the relevant time window into a single reference sentence.
All random seeds are fixed for reproducibility.

\paragraph{Scope distribution.}
With 55 streams, the benchmark generates approximately 2{,}500+ questions across the three scopes.
This distribution is driven by the six instant templates applied at four temporal anchor points per video, four recent templates at up to four time points, and seven historical templates at two to three time points.

\paragraph{Question templates.}
Six instant templates: \emph{``What is happening right now?''}, \emph{``Is anyone in the scene right now?''} (yes/no), \emph{``What can you see in the current frame?''}, \emph{``Describe the current scene.''}, \emph{``What objects are visible at this moment?''}, \emph{``Where does this scene take place?''}.

Four recent templates: \emph{``What just happened in the last few seconds?''}, \emph{``What was the person doing recently?''}, \emph{``Describe what occurred in the last moment.''}, \emph{``What changed recently in the scene?''}.

Seven historical templates: three open-ended (\emph{``What has happened throughout the stream?''}, \emph{``Summarize everything that has occurred so far.''}, \emph{``Describe the main events in the video so far.''}), three yes/no (\emph{``Did anything change earlier?''}, \emph{``Has there been more than one scene?''}, \emph{``Did the scene stay the same?''}), and one counting (\emph{``How many different scenes have appeared?''}).

\section{Demo Interface}
\label{sec:suppl_demo}

\begin{figure*}[t]
  \centering
  \includegraphics[width=\linewidth]{figures/demo_screenshot.png}
  \caption{\textbf{StreamMind live demo.} The interface shows a GTA~VI trailer playing in the video panel (left) with the SKM filmstrip and per-frame importance scores below it (Memory: 32/32). The chat panel (right) displays three queries spanning all temporal scopes: an \textsc{instant} query identifies the current scene, a \textsc{recent} query describes events from the last few seconds, and a \textsc{historical} query summarises all seven distinct scenes encountered so far. Each response includes its classified scope, the number of context frames used, and the end-to-end latency.}
  \label{fig:demo}
\end{figure*}

\Cref{fig:demo} shows the full demo interface. Users can stream from their webcam, upload a local video, or choose from pre-loaded sample clips including movie trailers and activity montages.
The video panel (left) shows the live feed with the SKM filmstrip; the chat panel (right) accepts free-form questions and returns scope-annotated answers.

\section{Evaluation Results and Reproducibility}
\label{sec:suppl_eval}

Full per-question evaluation results, ablation logs, and latency profiles are available in the project repository:

\begin{itemize}[itemsep=2pt,leftmargin=*]
  \item \textbf{LiveQA-Bench full results (A100):} \texttt{eval/results/streammind\_A100\_gpu\_results/} --- per-question predictions, accuracy breakdowns, and latency statistics across 55 video streams.
  \item \textbf{Ablation results:} \texttt{eval/results/streammind\_A100\_gpu\_results/ablation\_summary.json} --- accuracy and latency for each ablation configuration.
  \item \textbf{Latency profiling:} \texttt{eval/results/streammind\_A100\_gpu\_results/latency.json} --- per-component timing breakdown averaged over 10 profiling iterations.
  \item \textbf{Evaluation notebook:} \texttt{eval/StreamMind\_Eval.ipynb} --- Google Colab notebook for reproducing all results on a T4 or A100 GPU.
  \item \textbf{Project repository:} \url{https://github.com/palsure/StreamMind-LiveQA}
\end{itemize}

\noindent Results are organized as JSON files with per-question entries containing the predicted answer, ground-truth answer, scope label, accuracy score, and latency.
All evaluation uses fixed random seeds (seed=42) for reproducibility.
All models are frozen pre-trained checkpoints loaded from HuggingFace: CLIP ViT-B/32 (\texttt{openai/clip-vit-base-patch32}), BLIP captioning (\texttt{Salesforce/blip-image-captioning-base}), BLIP VQA (\texttt{Salesforce/blip-vqa-base}), and Flan-T5-base (\texttt{google/flan-t5-base}).
No training or fine-tuning is required.

{\small
\bibliographystyle{ieeenat_fullname}
\bibliography{main}
}

\end{document}
